\section{Unified Validation Across Angle, Frequency, and Polarization}

\subsection{Higher-order differentiation tests whether the same formalism handles sharper operator curvature}

If the proposed response-space formulation is genuinely more fundamental than a task-specific recipe, it should remain coherent when the target curvature changes without any change to the underlying inverse-design machinery. The fourth-order differentiation results provide that test. In the project archive, fourth-order candidates are scored using the same general logic adopted for the second-order benchmark, but with a target curve corresponding to a steeper dependence on transverse wavevector and with explicit tracking of center suppression, shape match, edge behavior, and effective bandwidth. The existence of viable high-scoring candidates across multiple wavelengths indicates that the present design strategy is not limited to one preselected transfer-function profile.

Scientifically, this matters because higher-order differentiation is not a trivial perturbation of the second-order case. Sharper target curvature places stronger demands on both the metasurface response and the ranking criteria. It is easier for a nominally promising design to match the center and fail at the edges, or to exhibit the right qualitative symmetry while missing the quantitative angular growth needed for useful operator fidelity. The fourth-order results therefore probe whether the response manifold being learned is broad enough to include multiple operator orders, rather than one favored response family. The answer suggested by the repository outputs is affirmative: the same candidate-generation and physical-reranking procedure remains meaningful when the operator order changes.

\subsection{Low-pass and high-pass targets test whether spectral-angular filtering can be absorbed into the same response language}

Angular low-pass and high-pass tasks probe a different axis of the unified formulation. Instead of changing the polynomial order of an operator-like target, they reshape the angular envelope itself. The repository implements both target families with score functions that search over reasonable effective angular widths and then evaluate center behavior, shape agreement, and bandwidth support. This is important because it converts what could have been presented as a new application into a simpler conceptual claim: angular filtering is another target geometry in response space.

The results support this reading. In the low-pass summaries, high-scoring candidates correspond to responses that preserve a strong central region while rejecting the outer angular band; in the high-pass summaries, the weighting is effectively reversed, favoring center suppression and strong outer-angle response. The same response input, candidate generation, and RCWA reranking procedure is used in both cases. No task-specific inverse architecture is introduced, and the geometry representation remains unchanged. These observations strengthen the central narrative of the manuscript. Spectral or angular filtering does not need to be framed as an ``extra task'' added onto a baseline differentiator recipe. It is already native to the response-space formulation because it is defined by a different target envelope over the same angle--wavelength coordinates.

\subsection{Polarization-independent and polarization-multiplexed designs test whether polarization enters as a first-class design axis}

The polarization-dependent tasks address perhaps the clearest possible challenge to the unified formulation. In many metasurface workflows, polarization is introduced late as an optional extension or as a separate design branch. Here it is built into the response tensor from the beginning. The repository includes both polarization-independent targets, which reward matched channel behavior, and polarization-multiplexed targets, which reward activity in one channel alongside suppression in the other. The score summaries explicitly report joint scores, channel-specific shape scores, passive-channel leakage, and channel values at large positive and negative angles. This is precisely the kind of bookkeeping one would expect if polarization were treated as a coequal part of the response object rather than as an add-on constraint.

The polarization-independent case asks whether the same design strategy can keep the two channels aligned over the target angular window. The polarization-multiplexed cases ask a harder question: can the same inverse formulation produce structures in which one channel remains functionally active while the other is driven toward silence? The answer emerging from the verified summaries is nuanced in a physically sensible way. Strong candidates do exist, but the gap between active-channel shape matching and passive-channel suppression reveals a genuine engineering tradeoff rather than a trivial design condition. That tradeoff is exactly what a physically meaningful formulation should expose. A response-space model that hid it behind a single averaged score would be less, not more, informative.

Taken together, the higher-order, filtering, and polarization-target results support the same conclusion. The contribution is not that one code path can be pointed at several unrelated operator tasks. The contribution is that those tasks can be represented, sampled, and verified within one angle--frequency--polarization design language. The unification is therefore conceptual before it is algorithmic: the response manifold is the right object to learn, and the cross-target results show that this object is rich enough to host multiple practically relevant Fourier-operator behaviors.
